\section{Human-in-the-loop 学习：从 TAMER 到 RLHF}

前面的案例主要依靠环境结果学习，本节加入人类作为额外反馈源。核心问题是怎样把稀疏、延迟且带主观噪声的评价转成可训练信号，同时避免 Agent 只迎合短期反馈而忽视长期任务结果。

Stone 特别回顾了 TAMER 系统（Teaching an Agent Manually via Evaluative Reinforcement）及其后续工作。它在 Tetris 上展示了 ``显式人类评价信号'' 如何显著加速早期学习。

\begin{figure}[H]
\centering
\includegraphics[width=0.95\textwidth]{diagrams/tamer-loop.png}
\caption{讲义整理图：TAMER 式人类评价、奖励建模与策略更新闭环}
\label{fig:tamer-loop}
\end{figure}

\begin{knowledgebox}{读图：人类反馈改变的是学习信号}
训练者对行为给正负评价，reward model 学习预测这种评价，policy 再据此更新。反馈不是直接遥控每一步，也不保证覆盖所有状态；因此后续环境奖励和自主探索仍不可缺。对应 LLM，偏好标签能塑造局部行为，但不能替代长程任务结果与安全约束。
\end{knowledgebox}

\begin{importantbox}{TAMER 的经验}
相对纯随机探索：
\begin{itemize}
    \item 学习早期收敛更快；
    \item 人类反馈能快速压制明显坏动作；
    \item 但最终上限可能受限于反馈噪声与覆盖度。
\end{itemize}
\end{importantbox}

\textbf{老师强调：}Stone 也指出后续组合路线应是 TAMER + RL。先用人类信号快速起步，再用自主 RL 继续提升上限。这几乎就是今天很多 LLM 后训练流程的结构映射：\textbf{human preference shaping + large-scale policy optimization}。

\begin{knowledgebox}{显式反馈与隐式反馈}
课程里还讨论了 implicit feedback（如乘客表情、紧张动作）：
\begin{itemize}
    \item 显式反馈：``good move / bad move''，高精度但高成本；
    \item 隐式反馈：无需额外标注流程，信号自然存在但噪声高。
\end{itemize}
这对应今天产品中 online telemetry 与 explicit rating 的组合。
\end{knowledgebox}

\begin{warningbox}{只靠 RLHF 不等于解决长期策略问题}
RLHF 能优化局部输出偏好，但在长链动作任务里，仍需配合环境反馈、过程奖励与安全约束，才能得到稳定 Agent 行为。
\end{warningbox}

\subsection{本章小结}
TAMER 的历史价值在于证明 ``人类反馈是可训练信号''。但完整 Agent 学习需要把人类反馈、环境反馈和自主探索放进同一系统。

